% -------------------------------------------------------------------------
% SUPERVISOR NOTE:
% This file is currently set up as a "Standalone Preview" so you can 
% compile it on Overleaf to verify the figure and caption look correct.
%
% WHEN INTEGRATING INTO YOUR MAIN PAPER:
% Copy ONLY the code between \begin{figure} and \end{figure} 
% and paste it into your 'main.tex' or 'theory_draft.tex'.
% -------------------------------------------------------------------------

\documentclass[10pt, a4paper]{article}

% --- REQUIRED PACKAGES (Copy these to your main paper's preamble if missing) ---
\usepackage{graphicx}  % Required for \includegraphics
\usepackage{caption}   % Better caption formatting
\usepackage{geometry}  % Sets page margins
\geometry{margin=1in}  % Gives the figure breathing room

\begin{document}

% -------------------------------------------------------------------------
% FIGURE 1: CONVERGENCE ANALYSIS
% Location: Results Section or Introduction
% -------------------------------------------------------------------------

\begin{figure}[t]
    \centering
    % NOTE: Ensure 'Figure_1_Convergence.pdf' is uploaded to your Overleaf project!
    \includegraphics[width=\columnwidth]{Figure_1_Convergence.pdf}
    \caption{
    \textbf{Convergence Comparison of Adaptive Controllers.} 
    \textit{(Point)} The proposed Physics-Informed Residual Learning (PIRL) architecture demonstrates a significant improvement in convergence speed compared to a standard model-free Neural Network. 
    \textit{(Evidence)} While the Pure NN (red) initializes with high variance ($RMSE > 0.5$ rad/s) and requires $>50$ epochs to approach the PID baseline, the PIRL agent (blue) leverages the Stribeck analytic prior to initialize at a near-converged state ($RMSE \approx 0.05$ rad/s). 
    \textit{(Explanation)} This effectively reduces the hypothesis space from the full friction manifold to a sparse residual term, allowing the controller to adapt to stiction non-linearities within the first 5 epochs.
    \textit{(Relay)} This rapid adaptation is critical for embedded deployment where training time is constrained by the 10ms control loop budget.
    }
    \label{fig:convergence}
\end{figure}

\end{document}